\chapter{第 2 章分层答案}
\label{app:ch02-solutions}

\section*{口述概念}
收敛序列必为 Cauchy 序列。Cauchy 序列只说明尾部元素彼此接近；只有空间完备时，它才保证收敛到空间内一点。$\mathbb{Q}$ 中逼近 $\sqrt2$ 的有理序列说明逆命题不能无条件使用。

\section*{笔试推导}
递推得 $d(x_{k+1},x_k)\le q^k d(x_1,x_0)$。对 $m>n$，
\[
d(x_m,x_n)\le\sum_{k=n}^{m-1}q^k d(x_1,x_0)
\le\frac{q^n}{1-q}d(x_1,x_0).
\]
这先证明序列是 Cauchy；完备性只在“存在 $x^*\in X$ 使 $x_n\to x^*$”这一步进入。随后由连续性得不动点，由 $q<1$ 得唯一性。

\section*{数值编程}
手算迭代依次为 $0.02,0.036,0.0488,0.05904,0.067232$，所以对固定点 $0.1$ 的误差是 $0.032768$，与先验界相同。对任意正整数 $n$，$x_n=2^{-1/n}$ 位于 $[0,1)$，故极限函数在该点为零，而 $x_n^n=1/2$；数值实验只复现这份解析见证。

\section*{研究判断}
连续小步长不是收敛证明。至少还要声明度量与参数尺度，证明或估计映射在不变集合上的压缩性质，说明状态空间完备，并检查固定点残差或目标函数误差。若映射并非压缩，应改用与算法匹配的收敛定理；有限 20 步只能作为诊断记录。
